\documentclass[10pt,a4paper]{article}
\usepackage[margin=15mm]{geometry}
\usepackage{fontspec}
\usepackage{xeCJK}
\usepackage{booktabs}
\usepackage{multirow}
\usepackage{siunitx}
\usepackage{caption}
\usepackage{pdflscape}
\usepackage{graphicx}
\usepackage{array}
\usepackage{xurl}
\usepackage[hidelinks]{hyperref}

\setmainfont{Times New Roman}
\setsansfont{Arial}
\setCJKmainfont{Microsoft YaHei}
\setCJKmonofont{Microsoft YaHei}
\sisetup{detect-weight=true,table-number-alignment=center}
\captionsetup{font=small,labelfont=bf,skip=4pt}
\setlength{\parindent}{0pt}
\setlength{\parskip}{0.45em}
\setlength{\tabcolsep}{4pt}
\pagestyle{plain}
\sloppy

\newcommand{\uprep}{$R_{\mathrm{up}}^{\mathrm{rep}}$}
\newcommand{\upwho}{$R_{\mathrm{up}}^{\mathrm{who}}$}
\newcommand{\onwho}{$R_{\mathrm{on}}^{\mathrm{who}}$}
\newcommand{\na}{---}
\newcommand{\zh}[1]{\par\noindent{\small #1}\par}

\begin{document}

\begin{center}
{\LARGE\bfseries Experiments / 实验}\\[4pt]
{\large BNPDFormer Main Results, Ablation Studies, and Baseline Comparisons}\\[2pt]
{\large BNPDFormer 主模型结果、消融实验与基线对比}
\end{center}

The experiments report BNPDFormer on the Min8 task at Start$\leq5$,
Start$\leq10$, and Start$\leq15$. Station identification, cause classification,
and remaining-time estimation are kept separate. Later sections remove the
graph, station-group embedding, delay pattern, or split will head, omit
process-heuristic labels and clustering, and compare B4 and B5 on the same
three start caps.
\zh{实验报告 BNPDFormer 在 Min8 任务、Start$\leq5/10/15$ 三档上的结果。工位识别、原因分类和剩余时间估计分开给出。后续各节分别去掉图结构、工位分组嵌入、延迟 pattern 或拆分的 will 头，去掉过程启发式标签和聚类，并在同样三个 Start 档位上与 B4、B5 对照。}

\section{Data and Protocol / 数据与实验协议}

\subsection{Task and Evaluation Protocol / 任务与评测协议}
The experiments evaluate station-level bottleneck-event prediction at three start
horizons, Start$\leq 5$, Start$\leq 10$, and Start$\leq 15$ min, while requiring
each target event to last at least eight minutes (Min8). The model observes a
30-min history and evaluates a 20-grid future label range at a one-minute
resolution. An \emph{ongoing} target is already active at the last history grid;
an \emph{upcoming} target starts after the observation boundary.
\zh{实验在三个起始时间范围内评估工位级瓶颈事件预测：Start$\leq5$、Start$\leq10$ 和 Start$\leq15$ 分钟，同时要求每个目标事件至少持续 8 分钟（Min8）。模型观察 30 分钟历史，并以 1 分钟为时间分辨率，在未来 20 个网格的标签范围内进行评测。若目标在历史序列最后一个网格已经发生，则记为进行中（ongoing）；若目标在观测边界之后才开始，则记为即将发生（upcoming）。}

The station-identification metrics are \texttt{will15}/\texttt{who} precision,
recall, and F1. The name \texttt{will15} is retained by the implementation but
refers to the active Start cap, rather than always denoting a fixed 15-min task.
The stricter report metric counts a true positive only when the station is
correct and $|\hat{s}-s|\leq 3$ min. Thus, \upwho measures station detection,
whereas \uprep additionally enforces start-time tolerance. Remaining-order-time
and event-duration errors are reported as mean absolute error (MAE; lower is
better). Process-cause performance is reported by accuracy, macro recall, and
per-class recall.
\zh{工位识别指标包括 \texttt{will15}/\texttt{who} 的精确率、召回率和 F1。代码沿用了 \texttt{will15} 这一名称，但该指标实际对应当前采用的 Start 上限，并不总是表示固定的 15 分钟任务。更严格的 report 指标只有在工位预测正确且 $|\hat{s}-s|\leq3$ 分钟时才计为真阳性。因此，\upwho 衡量是否识别出正确工位，而 \uprep 还额外要求起始时间满足容差。订单剩余时间和事件持续时间采用平均绝对误差（MAE，越低越好）；过程原因识别则报告准确率、宏平均召回率和逐类别召回率。}

All model and threshold decisions use validation data. The selected checkpoint
and threshold are then frozen for test evaluation. Results from different Start
caps are not pooled because their positive sets and upcoming denominators differ.
All reported runs use seed 42; consequently, no claim of statistical
significance is made.
\zh{所有模型选择和阈值选择均只使用验证集。确定 checkpoint 和阈值后，将二者冻结并用于测试集评测。由于不同 Start 上限对应的正例集合和 upcoming 分母不同，因此不合并不同 Start 档位的结果。所有实验均只使用 seed 42，所以本文不声称结果具有统计显著性。}

\subsection{Data Splits and Experiment Families / 数据划分与实验系列}
Two main-model settings share the BNPDFormer backbone. Separate report
fine-tuning selects each Start cap on its own. Sequential station fine-tuning
continues from the previous cap and selects on station F1. The full list of
changed settings is in Section~3.
\zh{两组主模型设置使用同一个 BNPDFormer 骨干。分档报告微调按 Start 档位分别选模。顺序工位微调沿用上一档 checkpoint，并按工位 F1 选模。具体改动列在第 3 节。}
\begin{center}
\captionof{table}{Experiment families and protocol boundaries; episode counts are train/validation/test. / 实验系列及协议边界；episode 数依次为训练/验证/测试。}
\small
\resizebox{\textwidth}{!}{
\begin{tabular}{llllll}
\toprule
Family & Dataset & Episodes & Main selection metric & Training budget & Test status \\
\midrule
Sequential station (retained)
 & \texttt{dense\_i1} & 140/28/36
 & validation \texttt{will15\_f1} & 100 epochs, patience 40 & evaluated \\
Separate report
 & \texttt{dense\_i1} & 140/28/36
 & validation \texttt{report\_f1} & 100 epochs, patience 100 & evaluated \\
No heuristic labels or clustering
 & \texttt{dense\_i1\_mix} & 166/37/45
 & validation \texttt{will15\_f1} & 100 epochs, patience 40 & evaluated \\
Structure ablations
 & \texttt{dense\_i1} & 140/28/36
 & validation \texttt{will15\_f1} & 50 epochs, from scratch & evaluated \\
B4 BTGCN / B5 BSTAN
 & fixed 208-episode set & 138/30/40
 & validation \texttt{will15\_f1} & 100 epochs, patience 40 & test unused \\
\bottomrule
\end{tabular}}
\end{center}

\section{Implementation Details / 实现细节}
The BNPDFormer encoder contains five blocks with an embedding dimension of 80
and an eight-dimensional Laplacian positional encoding. Graph, semantic, and
temporal attention use 2, 2, and 4 heads, respectively; these are separate
attention modules and should not be described as a single eight-head module.
The event head has hidden dimension 80. Training uses batch size 16, a 30-min
input history, one-minute windows, and a 20-grid future label range.
\zh{BNPDFormer 编码器包含 5 个块，嵌入维度为 80，并使用 8 维拉普拉斯位置编码。图注意力、语义注意力和时间注意力分别使用 2、2 和 4 个头；它们属于不同的注意力模块，因此不能表述为一个统一的“8 头注意力模块”。事件预测头的隐藏维度为 80。训练 batch size 为 16，输入历史长度为 30 分钟，时间窗口为 1 分钟，未来标签范围为 20 个网格。}

Optimization uses AdamW with a fine-tuning learning rate of $5\times10^{-5}$,
weight decay $0.05$, cosine learning-rate decay, and minimum learning rate
$10^{-6}$. The final runs allow at most 100 epochs and use early-stopping
patience 40. Event-positive windows are oversampled by a factor of 4. The model
is initialized from the earlier recall-oriented
\texttt{dense\_i1\_a1\_prefix8} checkpoint, unfreezes the last five encoder
blocks and occupancy decoder, and jointly optimizes event, occupancy,
remaining-time, and cause objectives. STGNPP and the legacy
\texttt{score}/\texttt{will}/\texttt{mark}/\texttt{tts} heads are disabled in
the retained configuration.
\zh{优化器采用 AdamW，微调学习率为 $5\times10^{-5}$，权重衰减为 $0.05$，使用余弦学习率衰减，最低学习率为 $10^{-6}$。最终实验最多训练 100 个 epoch，早停 patience 为 40；事件正例窗口按 4 倍过采样。模型从较早的召回导向 checkpoint \texttt{dense\_i1\_a1\_prefix8} 初始化，解冻最后 5 个编码器块和占用解码器，并联合优化事件、占用、剩余时间和原因任务。最终配置关闭 STGNPP，以及旧版 \texttt{score}/\texttt{will}/\texttt{mark}/\texttt{tts} 预测头。}

The retained setting is sequential station fine-tuning, configured by
\path{FactoryBN_dense_12_3_start{5,10,15}_min8_opt.json}.
Its runner loads the prefix8 checkpoint at Start$\leq5$, then warm-starts
Start$\leq10$ from the Start$\leq5$ checkpoint and Start$\leq15$ from
Start$\leq10$. Checkpoint selection uses validation station F1. The threshold
is chosen on validation and then frozen.
\zh{保留的设置是顺序工位微调，配置为 \texttt{FactoryBN\_dense\_12\_3\_start\{5,10,15\}\_min8\_opt.json}。训练脚本在 Start$\leq5$ 加载 prefix8 checkpoint，再让 Start$\leq10$ 从 Start$\leq5$ 热启动，Start$\leq15$ 从 Start$\leq10$ 热启动。选模使用验证集工位 F1。阈值在验证集上确定后冻结。}

\section{Main Results / 主实验结果}
\subsection{Separate Report Fine-tuning and Sequential Station Fine-tuning / 分档报告微调与顺序工位微调}
The two settings use the same BNPDFormer, the same \texttt{dense\_i1} split,
and the same Min8 task. Separate report fine-tuning is
\path{FactoryBN_dense_12_2_start{5,10,15}_min8_ep100.json}.
Each Start cap loads the prefix8 checkpoint on its own and keeps the parent
recipe: checkpoints are chosen by strict \texttt{report\_f1}, patience is 100,
and an ongoing station is lifted to the report threshold only when its
predicted duration is at least eight minutes. Sequential station fine-tuning
extends that Start$\leq5$ file. \path{run_start_horizon_min8_ep100.py} still
starts Start$\leq5$ from prefix8, then passes the previous best checkpoint to
the next cap. Selection changes to station \texttt{will15\_f1}.
\zh{两组设置使用同一个 BNPDFormer、同一份 \texttt{dense\_i1} 划分和同一个 Min8 任务。分档报告微调对应 \texttt{FactoryBN\_dense\_12\_2\_start\{5,10,15\}\_min8\_ep100.json}。每个 Start 档位各自加载 prefix8 checkpoint，并保留父配置：按严格 \texttt{report\_f1} 选模，patience 为 100；进行中工位只有在预测持续时间不少于 8 分钟时，才会被抬到报告阈值。顺序工位微调继承该 Start$\leq5$ 配置。\texttt{run\_start\_horizon\_min8\_ep100.py} 仍让 Start$\leq5$ 从 prefix8 开始，再把上一档最优 checkpoint 传给下一档。选模改为工位 \texttt{will15\_f1}。}

Both settings already share the validation gate $P\geq0.80$ and $R\geq0.70$,
split will heads, forced ongoing decisions, and the occupancy/upcoming union.
The parent upcoming positive weight is 8 and the far-start positive weight is
unset, so it stays 0. Events whose labeled start is within five minutes then
receive the shared near-start weight 6, which replaces the upcoming weight.
The weights 9/10/11 and the far-start loss boost of 2 therefore reach only
events that start later than five minutes. They do not change the Start$\leq5$
positive-class weights.
\zh{两边都已经带有验证集门槛 $P\geq0.80$、$R\geq0.70$、拆分的 will 头、进行中强制判定，以及占用与 upcoming 的并集。父配置的 upcoming 正例权重是 8，远起点正例权重没有设置，因此保持为 0。标注起点在 5 分钟以内的事件会改用两边共用的近起点权重 6，并覆盖 upcoming 权重。所以 9/10/11 这组权重和 2 倍远起点损失只作用在起点晚于 5 分钟的事件上，不改变 Start$\leq5$ 的正例权重。}

\begin{center}
\captionof{table}{Training and selection differences. Unlisted fields are inherited unchanged. / 训练与选模差异。未列出的字段沿用父配置，没有改动。}
\label{tab:recipe-diff}
\scriptsize
\resizebox{\textwidth}{!}{
\begin{tabular}{lll}
\toprule
Item & Separate report fine-tuning & Sequential station fine-tuning \\
\midrule
Initialization & each cap from prefix8 & Start$\leq5$ from prefix8, then $5\rightarrow10\rightarrow15$ \\
Checkpoint metric & \texttt{report\_f1} & \texttt{will15\_f1} \\
Patience & 100 & 40 \\
Upcoming / far positive weight & 8 / 0 at all caps & 9/9, 10/10, 11/11 \\
Far-start loss boost & 1 & 2, only if start $>5$ min \\
False-positive weight, hard multiplier & 2.4, 2.0 & 2.5, 2.2 \\
Ongoing duration requirement & on & off \\
Remaining-time loss & $w=0.35$, uniform & $w=0.5$, progress floor 0.25, power 1.5 \\
Threshold candidates & 0.55--0.98, uneven & 0.45--0.90 step 0.05, plus 0.94 \\
\bottomrule
\end{tabular}}
\end{center}

Table~\ref{tab:model-selection} compares the selected checkpoints. Separate
report fine-tuning is stronger at Start$\leq5$. Sequential station fine-tuning
is stronger at Start$\leq10$ and Start$\leq15$, and its mean validation F1 is
$0.845$ against $0.832$. It is retained because that validation margin holds
across the longer caps and the same selection rule is used for B4/B5.
\zh{表~\ref{tab:model-selection} 比较入选 checkpoint。分档报告微调在 Start$\leq5$ 更强。顺序工位微调在 Start$\leq10$ 和 Start$\leq15$ 更强，平均验证 F1 为 $0.845$，分档报告微调为 $0.832$。保留顺序工位微调，是因为较长预测范围上的验证优势成立，而且 B4/B5 使用同一条选模规则。}

\begin{center}
\captionof{table}{Selected checkpoints of the two fine-tuning settings. / 两种微调设置的入选 checkpoint。}
\label{tab:model-selection}
\scriptsize
\resizebox{\textwidth}{!}{
\begin{tabular}{llrrrrrrrrrr}
\toprule
Start & Family & Val $P$ & Val $R$ & Val F1 & Test $P$ & Test $R$ &
Test F1$_{\mathrm{who}}$ & Test F1$_{\mathrm{rep}}$ &
\uprep\ (test) & \onwho\ (test) & Best ep. \\
\midrule
$\leq5$  & Separate report & .865 & .872 & .869 & .832 & .808 & .820 & .819 & .542 & .923 & 33 \\
$\leq5$  & Sequential station & .837 & .891 & .863 & .794 & .817 & .805 & .804 & .539 & .937 & 39 \\
\midrule
$\leq10$ & Separate report & .805 & .836 & .820 & .750 & .763 & .757 & .751 & .550 & .923 & 83 \\
$\leq10$ & Sequential station & .831 & .840 & .835 & .806 & .761 & .783 & .778 & .528 & .936 & 41 \\
\midrule
$\leq15$ & Separate report & .814 & .803 & .808 & .772 & .716 & .743 & .741 & .489 & .925 & 37 \\
$\leq15$ & Sequential station & .841 & .835 & .838 & .806 & .742 & .772 & .768 & .534 & .925 & 42 \\
\bottomrule
\end{tabular}}
\end{center}

The separate-report Start$\leq5$ test station F1 is $0.820$, above the
sequential setting's $0.805$. At Start$\leq10$ and Start$\leq15$ the sequential
setting leads on both validation and test. The retained numbers in the
following tables are the sequential station fine-tuning runs.
\zh{分档报告微调在 Start$\leq5$ 的测试工位 F1 为 $0.820$，高于顺序工位微调的 $0.805$。Start$\leq10$ 和 Start$\leq15$ 上，顺序工位微调在验证集和测试集都更高。后面各表中的保留结果来自顺序工位微调。}

\begin{center}
\captionof{table}{Separate report fine-tuning station-timing test metrics. / 分档报告微调的工位时间测试指标。}
\scriptsize
\resizebox{\textwidth}{!}{
\begin{tabular}{lrrrrrrrr}
\toprule
Start & Dur. MAE & Start MAE & F1@1 & F1@2 & F1@3 &
\upwho & Threshold & $n_{\rm up}/n_{\rm on}/n_{\rm all}$ \\
\midrule
$\leq5$  & 1.93 & .13 & .801 & .814 & .819 & .545 & .90 & 336/766/1102 \\
$\leq10$ & 1.77 & .24 & .730 & .749 & .751 & .560 & .88 & 604/768/1372 \\
$\leq15$ & 1.88 & .28 & .706 & .735 & .741 & .489 & .94 & 708/770/1478 \\
\bottomrule
\end{tabular}}
\end{center}

\begin{center}
\captionof{table}{Separate report fine-tuning remaining-time results; MAE is in minutes. / 分档报告微调的剩余时间结果；MAE 单位为分钟。}
\scriptsize
\begin{tabular}{lrr}
\toprule
Start & Val remain MAE & Test remain MAE \\
\midrule
$\leq5$  & 12.03 & 10.64 \\
$\leq10$ & 10.86 &  9.61 \\
$\leq15$ & 11.77 & 10.20 \\
\bottomrule
\end{tabular}
\end{center}

\begin{center}
\captionof{table}{Separate report fine-tuning cause-classification results. / 分档报告微调的原因分类结果。}
\scriptsize
\begin{tabular}{lrrrrr}
\toprule
Start & Val accuracy & Val macro & Test accuracy & Test macro & Support \\
\midrule
$\leq5$  & .952 & .769 & .965 & .969 & 579 \\
$\leq10$ & .956 & .774 & .965 & .970 & 579 \\
$\leq15$ & .956 & .776 & .969 & .974 & 579 \\
\bottomrule
\end{tabular}
\end{center}

The train-derived majority-class accuracy is $0.409$. All three separate
report runs completed 100 epochs.
\zh{依据训练集得到的多数类基线准确率为 $0.409$。三组分档报告微调均训练满 100 个 epoch。}

\clearpage
\subsection{Final Main-Model Performance / 最终主模型表现}
\subsubsection{Station-Level Bottleneck Prediction / 工位级瓶颈预测}
\begin{center}
\captionof{table}{Sequential station fine-tuning results; validation is the selected epoch and test uses the frozen checkpoint/threshold. / 顺序工位微调结果；验证列来自入选 epoch，测试列使用冻结的 checkpoint 和阈值。}
\label{tab:main-event}
\scriptsize
\resizebox{\textwidth}{!}{
\begin{tabular}{llrrrrrrrrrr}
\toprule
Start & Split & $P$ & $R$ & F1$_{\mathrm{who}}$ & F1$_{\mathrm{rep}}$ &
\uprep & \onwho & $n_{\mathrm{up}}$ & $n_{\mathrm{on}}$ & $n_{\mathrm{all}}$ & Epoch \\
\midrule
$\leq5$  & Val  & .837 & .891 & .863 & .863 & .739 & .969 & \na & \na & \na & 39 \\
$\leq5$  & Test & .794 & .817 & \bfseries .805 & .804 & .539 & .937 & 336 & 766 & 1102 & 39 \\
\midrule
$\leq10$ & Val  & .831 & .840 & .835 & .835 & .711 & .960 & \na & \na & \na & 41 \\
$\leq10$ & Test & .806 & .761 & \bfseries .783 & .778 & .528 & .936 & 604 & 768 & 1372 & 41 \\
\midrule
$\leq15$ & Val  & .841 & .835 & .838 & .838 & .722 & .960 & \na & \na & \na & 42 \\
$\leq15$ & Test & .806 & .742 & \bfseries .772 & .768 & .534 & .925 & 708 & 770 & 1478 & 42 \\
\bottomrule
\end{tabular}}
\end{center}

The test station-level F1 decreases from $0.805$ to $0.772$ as the task includes
more distant upcoming events. Ongoing recall remains between $0.925$ and $0.937$,
whereas upcoming report recall remains around $0.53$. For Start$\leq15$, the
upcoming bucket recalls are $0.574$ for Start$\leq5$, $0.522$ for starts 6--10,
and $0.433$ for starts above 10 min. This identifies distant early warning,
rather than ongoing detection, as the principal limitation.
\zh{随着任务纳入更远的 upcoming 事件，测试集工位级 F1 从 $0.805$ 下降到 $0.772$。进行中召回率保持在 $0.925$--$0.937$，而 upcoming 严格报告召回率约为 $0.53$。在 Start$\leq15$ 设置下，Start$\leq5$、6--10 分钟以及大于 10 分钟三个 upcoming 桶的召回率分别为 $0.574$、$0.522$ 和 $0.433$。这说明主要限制是远距离提前预警能力，而不是进行中事件检测。}

\subsubsection{Remaining-Time Prediction / 剩余时间预测}
\begin{center}
\captionof{table}{Sequential station fine-tuning remaining-order-time results; MAE is in minutes. / 顺序工位微调的订单剩余时间结果；MAE 单位为分钟。}
\scriptsize
\begin{tabular}{llr}
\toprule
Start & Split & Remain MAE \\
\midrule
$\leq5$  & Val  & 12.3 \\
$\leq5$  & Test & 10.9 \\
$\leq10$ & Val  & 12.6 \\
$\leq10$ & Test & 11.2 \\
$\leq15$ & Val  & 11.0 \\
$\leq15$ & Test &  9.7 \\
\bottomrule
\end{tabular}
\end{center}

Test remaining-time MAE is $10.9$, $11.2$, and $9.7$ minutes at
Start$\leq5/10/15$. The error does not grow with the station-prediction
horizon: the Start$\leq15$ checkpoint is the most accurate of the three.
\zh{Start$\leq5/10/15$ 的测试集剩余时间 MAE 分别为 $10.9$、$11.2$ 和 $9.7$ 分钟。该误差并不随工位预测范围变长而上升：Start$\leq15$ 的 checkpoint 在三档中最准确。}

\subsubsection{Cause Classification / 瓶颈原因分类}
\begin{center}
\captionof{table}{Sequential station fine-tuning cause-classification results. / 顺序工位微调的瓶颈原因分类结果。}
\scriptsize
\begin{tabular}{llrrr}
\toprule
Start & Split & Cause accuracy & Cause macro recall & Cause support \\
\midrule
$\leq5$  & Val  & .956 & .963 & \na \\
$\leq5$  & Test & .965 & .967 & 579 \\
$\leq10$ & Val  & .958 & .968 & \na \\
$\leq10$ & Test & .962 & .965 & 579 \\
$\leq15$ & Val  & .967 & .978 & \na \\
$\leq15$ & Test & .974 & .978 & 579 \\
\bottomrule
\end{tabular}
\end{center}

\begin{center}
\captionof{table}{Sequential station fine-tuning per-class cause recall on test; the final two classes have no test support. / 顺序工位微调的测试集逐类别原因召回率；最后两类在测试集中无样本。}
\scriptsize
\begin{tabular}{lrrrrrr}
\toprule
Start & Transport delay & Material shortage & Starved upstream &
Queue buildup & Blocked downstream & High utilization \\
\midrule
$\leq5$  & .941 & .976 & .966 & .984 & no support & no support \\
$\leq10$ & .926 & .988 & .962 & .984 & no support & no support \\
$\leq15$ & .978 & .988 & .962 & .984 & no support & no support \\
\bottomrule
\end{tabular}
\end{center}

The cause head has four supported classes out of six configured classes.
The majority-class baseline is derived from train only and equals $0.409$.
Test accuracy stays between $0.962$ and $0.974$, well above that baseline.
The unsupported classes are retained in the table to make the support boundary
explicit rather than inflating a six-class claim.
\zh{原因识别头虽然配置了 6 个类别，但只有 4 个类别拥有有效样本。多数类基线只依据训练集确定，其准确率为 $0.409$。测试准确率保持在 $0.962$--$0.974$，明显高于该基线。表中仍保留两个无测试样本的类别，以明确支持度边界，避免将实际四类有效评测夸大为完整六类评测。}

\clearpage
\section{Ablation Studies / 消融实验}

\subsection{Ablation without Process-Heuristic Labels and Clustering / 未加入过程启发式标签与聚类的消融}
This ablation removes two components of the full method: the process-heuristic
labels and the clustering algorithm. It is trained from
\texttt{FactoryBN\_dense\_mix\_start\{5,10,15\}\_min8\_opt.json} on
\texttt{dense\_i1\_mix} (248 episodes, split 166/37/45). Selection uses
validation \texttt{will15\_f1}, and the warm-start order is
Start5$\rightarrow$Start10$\rightarrow$Start15. Start$\leq10/15$ use the
ungated validation best because the precision gate did not open.
\zh{该消融去掉完整方法中的两项组成：过程启发式标签和聚类算法。实验使用 \texttt{FactoryBN\_dense\_mix\_start\{5,10,15\}\_min8\_opt.json}，在 \texttt{dense\_i1\_mix}（248 个 episode，划分 166/37/45）上训练，以验证集 \texttt{will15\_f1} 选模，并按 Start5$\rightarrow$Start10$\rightarrow$Start15 顺序热启动。Start$\leq10/15$ 的精确率门槛未开启，因此采用无门槛验证集最优 checkpoint。}

\subsubsection{Station-Level Bottleneck Prediction / 工位级瓶颈预测}
\begin{center}
\captionof{table}{Station-level metrics without process-heuristic labels and clustering. / 未加入过程启发式标签与聚类时的工位级指标。}
\scriptsize
\resizebox{\textwidth}{!}{
\begin{tabular}{llrrrrrrrrrrrr}
\toprule
Start & Split & $P_{\mathrm{rep}}$ & $R_{\mathrm{rep}}$ & F1$_{\mathrm{rep}}$ &
$P_{\mathrm{who}}$ & $R_{\mathrm{who}}$ & F1$_{\mathrm{who}}$ &
\uprep & \upwho & \onwho & Support $n_{\rm up}/n_{\rm on}/n_{\rm all}$ &
Best/total ep. \\
\midrule
$\leq5$  & Val  & .822 & .718 & .766 & .824 & .720 & .769 & .349 & \na & .898 & \na & 1/41 \\
$\leq5$  & Test & .772 & .701 & .735 & .774 & .702 & .737 & .299 & .301 & .888 & 562/1216/1778 & 1/41 \\
\midrule
$\leq10$ & Val  & .728 & .682 & .704 & .742 & .695 & .718 & .460 & \na & .872 & \na & 58/98 \\
$\leq10$ & Test & .684 & .613 & .646 & .697 & .625 & .659 & .324 & .344 & .855 & 1000/1221/2221 & 58/98 \\
\midrule
$\leq15$ & Val  & .746 & .657 & .699 & .760 & .669 & .712 & .445 & \na & .868 & \na & 30/70 \\
$\leq15$ & Test & .697 & .586 & .637 & .704 & .592 & .643 & .320 & .325 & .847 & 1168/1222/2390 & 30/70 \\
\bottomrule
\end{tabular}}
\end{center}

Test report F1 falls to $.735/.646/.637$ at Start$\leq5/10/15$, which is
$.069/.132/.131$ below the retained model ($.804/.778/.768$). Upcoming report
recall also falls to $.299/.324/.320$. Station timing on the same checkpoints
is duration MAE $2.54/2.18/2.26$ min, start MAE $.13/.27/.22$ min, and
F1@1/2/3 equal to $.712/.729/.735$, $.626/.643/.646$, and $.623/.636/.637$.
\zh{Start$\leq5/10/15$ 的测试 report F1 降至 $.735/.646/.637$，比保留模型的 $.804/.778/.768$ 分别低 $.069/.132/.131$。Upcoming 严格召回率也降至 $.299/.324/.320$。同一 checkpoint 的工位时间指标为：持续时间 MAE $2.54/2.18/2.26$ 分钟，起点 MAE $.13/.27/.22$ 分钟，F1@1/2/3 分别为 $.712/.729/.735$、$.626/.643/.646$ 和 $.623/.636/.637$。}

\subsubsection{Remaining-Time Prediction / 剩余时间预测}
\begin{center}
\captionof{table}{Remaining-time results without process-heuristic labels and clustering; MAE is in minutes. / 未加入过程启发式标签与聚类时的剩余时间结果；MAE 单位为分钟。}
\scriptsize
\begin{tabular}{lrrrr}
\toprule
Start & Val MAE & Test MAE & Test MAE$_w$ & Test primary \\
\midrule
$\leq5$  & 12.96 & 11.85 & 9.71 & 9.56 \\
$\leq10$ & 11.58 & 10.95 & 9.12 & 9.77 \\
$\leq15$ & 11.64 & 10.92 & 9.03 & 9.57 \\
\bottomrule
\end{tabular}
\end{center}

MAE$_w$ is the progress-weighted global error and ``primary'' is the configured
middle-progress weighted error. Neither replaces the unweighted MAE. Test MAE
is higher than the retained model at Start$\leq5$ and Start$\leq15$
($11.85$ versus $10.9$, and $10.92$ versus $9.7$) and slightly lower at
Start$\leq10$ ($10.95$ versus $11.2$).
\zh{MAE$_w$ 是按生产进度加权的全局误差，“primary”是配置指定的中段进度加权误差；二者都不替代未加权 MAE。测试 MAE 在 Start$\leq5$ 和 Start$\leq15$ 高于保留模型（$11.85$ 对 $10.9$，$10.92$ 对 $9.7$），在 Start$\leq10$ 略低（$10.95$ 对 $11.2$）。}

\subsubsection{Cause Classification / 瓶颈原因分类}
\begin{center}
\captionof{table}{Cause-classification results without process-heuristic labels and clustering. / 未加入过程启发式标签与聚类时的原因分类结果。}
\scriptsize
\begin{tabular}{lrrrrr}
\toprule
Start & Val accuracy & Val macro & Test accuracy & Test macro & Support \\
\midrule
$\leq5$  & .931 & .942 & .944 & .964 & 964 \\
$\leq10$ & .928 & .937 & .924 & .940 & 964 \\
$\leq15$ & .938 & .939 & .940 & .956 & 964 \\
\bottomrule
\end{tabular}
\end{center}

\begin{center}
\captionof{table}{Per-class cause recall on test without process-heuristic labels and clustering. / 未加入过程启发式标签与聚类时的测试集逐类别原因召回率。}
\scriptsize
\begin{tabular}{lrrrr}
\toprule
Start & Transport delay & Material shortage & Starved upstream & Queue buildup \\
\midrule
$\leq5$  & .962 & .989 & .906 & 1.000 \\
$\leq10$ & .928 & .956 & .893 & .984 \\
$\leq15$ & .943 & .978 & .910 & .995 \\
\bottomrule
\end{tabular}
\end{center}

The train-derived majority baseline on this split is $0.495$. Test accuracy
($.944/.924/.940$) remains far above that baseline, but it is below the retained
model ($.965/.962/.974$). Together with the station-level drop, the arm shows
that leaving out process-heuristic labels and clustering weakens both event
detection and cause classification.
\zh{该划分上由训练集得到的多数类基线为 $0.495$。测试准确率（$.944/.924/.940$）仍远高于该基线，但低于保留模型（$.965/.962/.974$）。结合工位级下降，该消融表明：不加入过程启发式标签和聚类，会同时削弱事件检测和原因分类。}

\clearpage
\subsection{Model-Structure Ablations / 模型结构消融}
The \texttt{nograph} arm sets \texttt{graph\_mode=identity}, removing graph
edges. The \texttt{nogroup} arm removes station-group/semantic-group embedding.
\texttt{nopattern} disables the delayed pattern injection in the geographic
attention keys. \texttt{nosplit} uses one shared event-will head for upcoming
and ongoing bottlenecks (\texttt{split\_will\_heads=false}); decode-time
forcing of ongoing events is left unchanged. All four arms use the same
\texttt{dense\_i1} split as the final model. Start$\leq5$ is trained from
scratch for at most 50 epochs with an empty \texttt{init\_ckpt}; Start$\leq10$
and Start$\leq15$ warm-start from the previous arm checkpoint. This budget
differs from the warm-started 100-epoch full model and is reported as a
limitation of causal interpretation.
\zh{\texttt{nograph} 将 \texttt{graph\_mode} 设为 \texttt{identity}，从而移除图边；\texttt{nogroup} 移除工位分组/语义分组嵌入；\texttt{nopattern} 关闭地理注意力键中的延迟 pattern 注入；\texttt{nosplit} 让 upcoming 与 ongoing 共用一个事件 will 头（\texttt{split\_will\_heads=false}），解码时对 ongoing 的强制保持不变。四组消融都使用与最终模型相同的 \texttt{dense\_i1} 划分。Start$\leq5$ 以空 \texttt{init\_ckpt} 从头训练，上限 50 个 epoch；Start$\leq10$ 和 Start$\leq15$ 从上一级 checkpoint 热启动。该预算不同于热启动且最多训练 100 个 epoch 的完整模型，因此这是解释消融因果效果时必须保留的限制。}

\subsubsection{Station-Level Bottleneck Prediction / 工位级瓶颈预测}
\begin{center}
\captionof{table}{Structure-ablation station metrics at Start$\leq5$; Report F1 requires a correct station and a start error within 3 min. / Start$\leq5$ 结构消融工位指标；Report F1 要求工位正确且起点误差不超过 3 分钟。}
\scriptsize
\resizebox{\textwidth}{!}{
\begin{tabular}{llrrrrrrrrrr}
\toprule
Arm & Split & $P$ & $R$ & F1$_{\mathrm{who}}$ & F1$_{\mathrm{rep}}$ &
\uprep & \onwho & $n_{\rm up}$ & $n_{\rm on}$ & $n_{\rm all}$ & Best ep. \\
\midrule
Full       & Val  & .837 & .891 & .863 & .863 & .739 & .969 & \na & \na & \na & 39 \\
Full       & Test & .794 & .817 & .805 & .804 & .539 & .937 & 336 & 766 & 1102 & 39 \\
No graph   & Val  & .823 & .658 & .731 & .731 & .115 & .937 & \na & \na & \na & 19 \\
No graph   & Test & .832 & .678 & .747 & .746 & .101 & .930 & 336 & 766 & 1102 & 19 \\
No group   & Val  & .782 & .765 & .774 & .774 & .383 & .962 & \na & \na & \na & 28 \\
No group   & Test & .766 & .756 & .761 & .753 & .307 & .941 & 336 & 766 & 1102 & 28 \\
No pattern & Val  & .789 & .671 & .725 & .725 & .166 & .930 & \na & \na & \na & 15 \\
No pattern & Test & .799 & .698 & .748 & .745 & .158 & .935 & 336 & 766 & 1102 & 15 \\
No split   & Val  & .784 & .670 & .724 & .723 & .149 & .937 & \na & \na & \na & 15 \\
No split   & Test & .750 & .706 & .729 & .727 & .185 & .935 & 336 & 766 & 1102 & 15 \\
\bottomrule
\end{tabular}}
\end{center}

\begin{center}
\captionof{table}{Structure-ablation station metrics at Start$\leq10$. / Start$\leq10$ 结构消融工位指标。}
\scriptsize
\resizebox{\textwidth}{!}{
\begin{tabular}{llrrrrrrrrrr}
\toprule
Arm & Split & $P$ & $R$ & F1$_{\mathrm{who}}$ & F1$_{\mathrm{rep}}$ &
\uprep & \onwho & $n_{\rm up}$ & $n_{\rm on}$ & $n_{\rm all}$ & Best ep. \\
\midrule
Full       & Val  & .831 & .840 & .835 & .835 & .711 & .960 & \na & \na & \na & 41 \\
Full       & Test & .806 & .761 & .783 & .778 & .528 & .936 & 604 & 768 & 1372 & 41 \\
No graph   & Val  & .807 & .486 & .606 & .606 & .049 & .894 & \na & \na & \na & 9 \\
No graph   & Test & .821 & .539 & .651 & .649 & .065 & .909 & 604 & 768 & 1372 & 9 \\
No group   & Val  & .663 & .790 & .721 & .721 & .636 & .934 & \na & \na & \na & 44 \\
No group   & Test & .605 & .745 & .668 & .651 & .445 & .947 & 604 & 768 & 1372 & 44 \\
No pattern & Val  & .598 & .659 & .643 & .627 & .360 & .939 & \na & \na & \na & 28 \\
No pattern & Test & .555 & .625 & .607 & .588 & .235 & .934 & 604 & 768 & 1372 & 28 \\
No split   & Val  & .596 & .614 & .629 & .605 & .254 & .951 & \na & \na & \na & 33 \\
No split   & Test & .546 & .618 & .599 & .580 & .204 & .947 & 604 & 768 & 1372 & 33 \\
\bottomrule
\end{tabular}}
\end{center}

\begin{center}
\captionof{table}{Structure-ablation station metrics at Start$\leq15$. / Start$\leq15$ 结构消融工位指标。}
\scriptsize
\resizebox{\textwidth}{!}{
\begin{tabular}{llrrrrrrrrrr}
\toprule
Arm & Split & $P$ & $R$ & F1$_{\mathrm{who}}$ & F1$_{\mathrm{rep}}$ &
\uprep & \onwho & $n_{\rm up}$ & $n_{\rm on}$ & $n_{\rm all}$ & Best ep. \\
\midrule
Full       & Val  & .841 & .835 & .838 & .838 & .722 & .960 & \na & \na & \na & 42 \\
Full       & Test & .806 & .742 & .772 & .768 & .534 & .925 & 708 & 770 & 1478 & 42 \\
No graph   & Val  & .819 & .435 & .568 & .568 & .028 & .883 & \na & \na & \na & 11 \\
No graph   & Test & .851 & .482 & .616 & .612 & .037 & .887 & 708 & 770 & 1478 & 11 \\
No group   & Val  & .620 & .760 & .683 & .683 & .609 & .927 & \na & \na & \na & 50 \\
No group   & Test & .567 & .690 & .623 & .604 & .395 & .922 & 708 & 770 & 1478 & 50 \\
No pattern & Val  & .631 & .769 & .705 & .693 & .600 & .956 & \na & \na & \na & 42 \\
No pattern & Test & .539 & .681 & .616 & .601 & .397 & .947 & 708 & 770 & 1478 & 42 \\
No split   & Val  & .627 & .690 & .672 & .657 & .453 & .951 & \na & \na & \na & 47 \\
No split   & Test & .539 & .641 & .607 & .586 & .309 & .951 & 708 & 770 & 1478 & 47 \\
\bottomrule
\end{tabular}}
\end{center}

\subsubsection{Remaining-Time Prediction / 剩余时间预测}
\begin{center}
\captionof{table}{Remaining-time MAE (minutes) for the structure ablations. / 结构消融的剩余时间 MAE（分钟）。}
\scriptsize
\begin{tabular}{lllrr}
\toprule
Arm & Start & Val MAE & Test MAE \\
\midrule
Full       & $\leq5$  & 12.3 & 10.9 \\
No graph   & $\leq5$  & 18.2 & 17.0 \\
No group   & $\leq5$  & 13.2 & 12.7 \\
No pattern & $\leq5$  & 20.1 & 18.6 \\
No split   & $\leq5$  & 20.8 & 18.9 \\
\midrule
Full       & $\leq10$ & 12.6 & 11.2 \\
No graph   & $\leq10$ & 24.1 & 23.3 \\
No group   & $\leq10$ & 12.5 & 11.6 \\
No pattern & $\leq10$ & 14.8 & 13.4 \\
No split   & $\leq10$ & 15.6 & 14.4 \\
\midrule
Full       & $\leq15$ & 11.0 &  9.7 \\
No graph   & $\leq15$ & 21.4 & 20.8 \\
No group   & $\leq15$ & 11.1 & 10.6 \\
No pattern & $\leq15$ & 12.5 & 11.1 \\
No split   & $\leq15$ & 15.2 & 13.8 \\
\bottomrule
\end{tabular}
\end{center}

\subsubsection{Cause Classification / 瓶颈原因分类}
\begin{center}
\captionof{table}{Cause-classification metrics for the structure ablations. / 结构消融的原因分类指标。}
\scriptsize
\resizebox{\textwidth}{!}{
\begin{tabular}{lllrrrrr}
\toprule
Arm & Start & Split & Accuracy & Macro recall & Majority baseline & Support \\
\midrule
Full       & $\leq5$  & Val  & .956 & .963 & .409 & \na \\
Full       & $\leq5$  & Test & .965 & .967 & .409 & 579 \\
No graph   & $\leq5$  & Val  & .866 & .886 & .409 & \na \\
No graph   & $\leq5$  & Test & .876 & .889 & .409 & 579 \\
No group   & $\leq5$  & Val  & .932 & .952 & .409 & \na \\
No group   & $\leq5$  & Test & .943 & .955 & .409 & 579 \\
No pattern & $\leq5$  & Val  & .726 & .725 & .409 & \na \\
No pattern & $\leq5$  & Test & .732 & .734 & .409 & 579 \\
No split   & $\leq5$  & Val  & .853 & .874 & .409 & \na \\
No split   & $\leq5$  & Test & .853 & .874 & .409 & 579 \\
\midrule
Full       & $\leq10$ & Val  & .958 & .968 & .409 & \na \\
Full       & $\leq10$ & Test & .962 & .965 & .409 & 579 \\
No graph   & $\leq10$ & Val  & .741 & .755 & .409 & \na \\
No graph   & $\leq10$ & Test & .734 & .746 & .409 & 579 \\
No group   & $\leq10$ & Val  & .928 & .944 & .409 & \na \\
No group   & $\leq10$ & Test & .938 & .953 & .409 & 579 \\
No pattern & $\leq10$ & Val  & .807 & .823 & .409 & \na \\
No pattern & $\leq10$ & Test & .793 & .805 & .409 & 579 \\
No split   & $\leq10$ & Val  & .888 & .904 & .409 & \na \\
No split   & $\leq10$ & Test & .869 & .884 & .409 & 579 \\
\midrule
Full       & $\leq15$ & Val  & .967 & .978 & .409 & \na \\
Full       & $\leq15$ & Test & .974 & .978 & .409 & 579 \\
No graph   & $\leq15$ & Val  & .796 & .790 & .409 & \na \\
No graph   & $\leq15$ & Test & .794 & .787 & .409 & 579 \\
No group   & $\leq15$ & Val  & .941 & .950 & .409 & \na \\
No group   & $\leq15$ & Test & .938 & .949 & .409 & 579 \\
No pattern & $\leq15$ & Val  & .851 & .862 & .409 & \na \\
No pattern & $\leq15$ & Test & .843 & .850 & .409 & 579 \\
No split   & $\leq15$ & Val  & .897 & .915 & .409 & \na \\
No split   & $\leq15$ & Test & .869 & .890 & .409 & 579 \\
\bottomrule
\end{tabular}}
\end{center}

\begin{landscape}
\begin{center}
\captionof{table}{Complete per-class cause recall for the full model and structure ablations on test. / 完整模型与结构消融在测试集上的逐类别原因召回率。}
\scriptsize
\begin{tabular}{llrrrrrr}
\toprule
Arm & Start & Transport delay & Material shortage & Starved upstream &
Queue buildup & Blocked downstream & High utilization \\
\midrule
Full     & $\leq5$  & .941 & .976 & .966 & .984 & no support & no support \\
No graph & $\leq5$  & .971 & 1.000 & .852 & .734 & no support & no support \\
No group & $\leq5$  & .956 & 1.000 & .911 & .952 & no support & no support \\
No pattern & $\leq5$  & .993 & 1.000 & .797 & .145 & no support & no support \\
No split   & $\leq5$  & .934 & 1.000 & .806 & .758 & no support & no support \\
\midrule
Full     & $\leq10$ & .926 & .988 & .962 & .984 & no support & no support \\
No graph & $\leq10$ & .993 & .963 & .738 & .290 & no support & no support \\
No group & $\leq10$ & .956 & 1.000 & .890 & .968 & no support & no support \\
No pattern & $\leq10$ & .941 & 1.000 & .802 & .476 & no support & no support \\
No split   & $\leq10$ & .831 & 1.000 & .857 & .847 & no support & no support \\
\midrule
Full     & $\leq15$ & .978 & .988 & .962 & .984 & no support & no support \\
No graph & $\leq15$ & .993 & .988 & .878 & .290 & no support & no support \\
No group & $\leq15$ & .926 & 1.000 & .911 & .960 & no support & no support \\
No pattern & $\leq15$ & .971 & 1.000 & .857 & .573 & no support & no support \\
No split   & $\leq15$ & .875 & 1.000 & .823 & .863 & no support & no support \\
\bottomrule
\end{tabular}
\end{center}
\end{landscape}

Removing graph structure causes the largest and most consistent degradation,
especially for upcoming recall, remaining-time MAE, and queue-buildup recall.
Removing group embeddings and delayed pattern injection also lowers station F1
at every horizon. The \texttt{nopattern} test report-F1 drops are
$.059/.190/.167$, and its queue-buildup recall is substantially lower.
Sharing one will head (\texttt{nosplit}) drops test report F1 by
$.077/.198/.182$. Upcoming report recall falls to $.185/.204/.309$, while
ongoing recall stays above $.93$. Remaining-time MAE rises to
$18.9/14.4/13.8$ minutes. Cause accuracy stays high ($.853/.869/.869$) because
the cause head is unchanged; the damage is concentrated in station detection.
The \texttt{nogroup} validation precision values ($.782/.663/.620$) fail the
$P\geq.80$ gate, and the gate also never opens for \texttt{nopattern} or
\texttt{nosplit}; those rows use the ungated validation best. The
\texttt{nogroup} Start$\leq15$ checkpoint is the final epoch (50), so it did
not plateau within the assigned budget.
\zh{移除图结构造成最大且最一致的性能下降，尤其体现在 upcoming 召回率、剩余时间 MAE 和 queue-buildup 类召回率上。移除工位分组嵌入和延迟 pattern 注入也会降低三个 Start 档位的工位 F1。\texttt{nopattern} 的测试 report-F1 分别下降 $.059/.190/.167$，且 queue-buildup 类召回率明显降低。共用一个 will 头（\texttt{nosplit}）使测试 report F1 下降 $.077/.198/.182$，upcoming 严格召回率降至 $.185/.204/.309$，而 ongoing 召回率仍高于 $.93$。剩余时间 MAE 升至 $18.9/14.4/13.8$ 分钟。原因准确率仍较高（$.853/.869/.869$），因为原因头没有改动，损失集中在工位检测。\texttt{nogroup} 的验证集精确率 $.782/.663/.620$ 未达到 $P\geq.80$；\texttt{nopattern} 和 \texttt{nosplit} 的精确率门槛也从未开启，这些行采用无门槛验证集最优值。\texttt{nogroup} 的 Start$\leq15$ checkpoint 位于最后一个 epoch（50），说明该实验在给定预算内尚未进入平台期。}

The planned \texttt{nosem} arm
(\texttt{ablation\_nosem\_start5\_seed42}) aborted before epoch 1 because a
convolution weight had shape \texttt{[0,80,1,1]}; consequently, it has no valid
metrics. This failed arm is reported rather than silently excluded.
\zh{计划中的 \texttt{nosem} 实验（\texttt{ablation\_nosem\_start5\_seed42}）在第 1 个 epoch 之前即终止，原因是卷积权重形状为 \texttt{[0,80,1,1]}，因此没有有效指标。本文明确报告该失败实验，而不是将其静默删除。}

\clearpage
\section{Baseline Comparison / 基线对比}

\subsection{Comparable Setup and Limitations / 对齐设置与局限}
B4 is BTGCN with a GCN+GRU128 backbone; B5 is BSTAN with scalar
GAT+GRU128. Each baseline retains its own near-history input and parent
checkpoint and does not receive BNPDFormer-specific branches. Training follows
Start5$\rightarrow$Start10$\rightarrow$Start15. The six runs use the same
Start/Min8 semantics, 20 future grids, threshold candidates, 3-min start
tolerance, and validation checkpoint rule as the final BNPDFormer protocol.
All six use the fallback F1 rule because none meets both $P\geq.80$ and
$R\geq.70$.
\zh{B4 为采用 GCN+GRU128 骨干的 BTGCN，B5 为采用标量 GAT+GRU128 骨干的 BSTAN。每个基线保留自身的 near 历史输入和父 checkpoint，不引入 BNPDFormer 专有分支。训练顺序为 Start5$\rightarrow$Start10$\rightarrow$Start15。六组实验与最终 BNPDFormer 采用相同的 Start/Min8 语义、20 个未来网格、阈值候选集合、3 分钟起点容差以及验证集选模规则。由于六组都未同时满足 $P\geq.80$ 和 $R\geq.70$，最终均采用 F1 fallback 规则。}

The baseline split is 138/30/40 rather than 140/28/36, and the 40 baseline test
episodes have not been evaluated. Therefore, Tables~\ref{tab:b45-val} and
\ref{tab:b45-train} are complete validation/training reports, but they are not
test-set evidence that BNPDFormer outperforms B4/B5.
\zh{基线的数据划分为 138/30/40，而不是主模型的 140/28/36，并且基线的 40 个测试 episode 尚未参与评测。因此，表~\ref{tab:b45-val} 和表~\ref{tab:b45-train} 是完整的验证集/训练集报告，但不能作为 BNPDFormer 在测试集上优于 B4/B5 的证据。}

\begin{landscape}
\begin{center}
\captionof{table}{All official B4/B5 validation metrics; hit counts are over upcoming support. / B4/B5 全部正式验证指标；命中数均以 upcoming 支持数为分母。}
\label{tab:b45-val}
\scriptsize
\resizebox{0.97\linewidth}{!}{
\begin{tabular}{llrrrrllllrrr}
\toprule
Model & Start & $P$ & $R$ & F1 & \uprep & \upwho & \onwho &
Upcoming strict & Upcoming station & Best/total ep. & Threshold \\
\midrule
B4 & $\leq5$  & .8048 & .5655 & .6643 & .0742 & .0742 & .8078 &
32/431 & 32/431 & 1/41 & .50 \\
B5 & $\leq5$  & .8133 & .5709 & .6709 & .0812 & .0812 & .8124 &
35/431 & 35/431 & 2/42 & .70 \\
B4 & $\leq10$ & .8152 & .3829 & .5211 & .0353 & .0353 & .6833 &
27/764 & 27/764 & 3/43 & .90 \\
B5 & $\leq10$ & .8118 & .4163 & .5503 & .0314 & .0340 & .7466 &
24/764 & 26/764 & 2/42 & .85 \\
B4 & $\leq15$ & .6684 & .3718 & .4778 & .0280 & .0325 & .7133 &
25/892 & 29/892 & 5/45 & .94 \\
B5 & $\leq15$ & .8036 & .3543 & .4918 & .0291 & .0336 & .6772 &
26/892 & 30/892 & 2/42 & .90 \\
\bottomrule
\end{tabular}}
\end{center}
\end{landscape}

\begin{center}
\captionof{table}{B4/B5 training station metrics. / B4/B5 训练集工位指标。}
\label{tab:b45-train}
\scriptsize
\begin{tabular}{llrlr}
\toprule
Model & Start & Train F1 & Train upcoming strict hits & Train \uprep \\
\midrule
B4 & $\leq5$  & .6918 & 98/1753   & .0559 \\
B5 & $\leq5$  & .7216 & 170/1753  & .0970 \\
B4 & $\leq10$ & .6349 & 324/3109  & .1042 \\
B5 & $\leq10$ & .6642 & 295/3109  & .0949 \\
B4 & $\leq15$ & .8816 & 2657/3624 & .7332 \\
B5 & $\leq15$ & .6434 & 477/3624  & .1316 \\
\bottomrule
\end{tabular}
\end{center}

\begin{center}
\captionof{table}{B4/B5 validation remaining-time MAE; values are middle-progress-weighted and kept in window units. / B4/B5 验证集剩余时间 MAE；数值为中段进度加权，单位保留为窗口。}
\scriptsize
\begin{tabular}{llr}
\toprule
Model & Start & Val remain MAE \\
\midrule
B4 & $\leq5$  & 13.3421 \\
B5 & $\leq5$  & 11.2884 \\
B4 & $\leq10$ & 16.2363 \\
B5 & $\leq10$ & 10.6643 \\
B4 & $\leq15$ & 11.7820 \\
B5 & $\leq15$ &  9.6109 \\
\bottomrule
\end{tabular}
\end{center}

\begin{center}
\captionof{table}{B4/B5 validation cause-classification metrics. / B4/B5 验证集原因分类指标。}
\scriptsize
\begin{tabular}{llrr}
\toprule
Model & Start & Val cause accuracy & Val cause macro recall \\
\midrule
B4 & $\leq5$  & .8876 & .8631 \\
B5 & $\leq5$  & .9575 & .9538 \\
B4 & $\leq10$ & .9001 & .8750 \\
B5 & $\leq10$ & .9572 & .9481 \\
B4 & $\leq15$ & .9040 & .8837 \\
B5 & $\leq15$ & .9518 & .9362 \\
\bottomrule
\end{tabular}
\end{center}

B5 has higher validation F1 than B4 by $0.0066$, $0.0292$, and $0.0140$ at the
three horizons and lower remaining-time MAE in every pair. It does not dominate
strict upcoming recall: at Start$\leq10$, B5 obtains $0.0314$ versus B4's
$0.0353$. Both baselines remain weak on upcoming events. The most severe
generalization gap occurs for B4 Start$\leq15$, where train F1/upcoming strict
recall are $0.8816/0.7332$, but validation F1/upcoming strict recall fall to
$0.4778/0.0280$.
\zh{B5 在三个 Start 档位上的验证 F1 分别比 B4 高 $0.0066$、$0.0292$ 和 $0.0140$，且每一档的剩余时间 MAE 都更低。但 B5 并未在严格 upcoming 召回率上全面占优：Start$\leq10$ 时，B5 为 $0.0314$，低于 B4 的 $0.0353$。两个基线在 upcoming 事件上都较弱。最严重的泛化差距出现在 B4 Start$\leq15$：训练 F1/严格 upcoming 召回率为 $0.8816/0.7332$，验证集则降至 $0.4778/0.0280$。}

\subsection{Overall Discussion and Validity Boundaries / 总体讨论与有效性边界}
The retained sequential station fine-tuning reaches test station-level F1 of
$0.805/0.783/0.772$ and strict upcoming recall of
$0.539/0.528/0.534$ for Start$\leq5/10/15$. The graph, station-group, pattern,
and shared-head ablations all reduce station F1. Graph removal nearly
eliminates upcoming recall at long horizons, and the shared will head cuts
upcoming report recall to $0.185/0.204/0.309$. Leaving out process-heuristic
labels and clustering lowers test report F1 to $0.735/0.646/0.637$. Cause
classification and remaining-time estimation are reported separately from
station identification; neither is folded into the station F1.
\zh{保留的顺序工位微调在 Start$\leq5/10/15$ 三档上的测试集工位级 F1 分别为 $0.805/0.783/0.772$，严格 upcoming 召回率分别为 $0.539/0.528/0.534$。移除图结构、工位分组、延迟 pattern，或让 upcoming 与 ongoing 共用 will 头，都会降低工位 F1。移除图结构几乎使长预测范围的 upcoming 召回能力消失；共用 will 头则把 upcoming 严格召回率降到 $0.185/0.204/0.309$。不加入过程启发式标签和聚类时，测试 report F1 降至 $0.735/0.646/0.637$。原因分类和剩余时间估计与工位识别分开报告，不并入工位 F1。}

Four validity limits must accompany these results. First, all runs use one seed,
so uncertainty and significance are unknown. Second, B4/B5 currently have no
test result and use a different episode split; their validation scores cannot be
presented as a final test ranking. Third, the structure ablations use a smaller,
from-scratch budget than the warm-started full model. Fourth, remaining-time
metrics are not uniformly comparable: BNPDFormer reports minute-scale
\texttt{remain\_len\_mae}, while B4/B5 retain middle-progress-weighted values in
window units. These boundaries are preserved in the tables rather than being
hidden by a single pooled score.
\zh{这些结果必须同时说明四项有效性限制。第一，所有实验只使用一个随机种子，因此无法估计不确定性或统计显著性。第二，B4/B5 尚无测试集结果，且采用不同的 episode 划分，所以不能将其验证分数作为最终测试排名。第三，结构消融采用比完整模型更小的、从头训练的预算。第四，剩余时间指标并非完全可比：BNPDFormer 报告以分钟为单位的 \texttt{remain\_len\_mae}，B4/B5 则保留以窗口为单位的中段进度加权值。本文在各表中明确保留这些边界，而不是用一个混合总分掩盖差异。}

\end{document}
